%bismi ALLAH ARRAHMAN ARRAHIM









































































\documentclass[11pt,a4paper]{report}

% ============================================================
%  THEORIE DU MAILLAGE CONTEXTUEL ET DE LA NAVIGATION CAUSALE
% ============================================================

\usepackage[utf8]{inputenc}
\usepackage[T1]{fontenc}
\usepackage[french]{babel}
\usepackage{lmodern}
\usepackage{geometry}
\usepackage{amsmath,amssymb,amsthm,mathtools}
\usepackage{enumitem}
\usepackage{hyperref}
\usepackage{microtype}

\geometry{margin=2.5cm}
\setlength{\parskip}{0.45em}
\setlength{\parindent}{0pt}

\hypersetup{
  colorlinks=true,
  linkcolor=blue,
  urlcolor=blue,
  pdftitle={Theorie du maillage contextuel et de la navigation causale},
  pdfauthor={}
}

% ----------------------------
% Environnements mathematiques
% ----------------------------

\theoremstyle{definition}
\newtheorem{definition}{Définition}[chapter]
\newtheorem{postulat}[definition]{Postulat}
\newtheorem{hypothese}[definition]{Hypothèse}
\newtheorem{convention}[definition]{Convention}

\theoremstyle{plain}
\newtheorem{proposition}[definition]{Proposition}
\newtheorem{theoreme}[definition]{Théorème}
\newtheorem{lemme}[definition]{Lemme}
\newtheorem{corollaire}[definition]{Corollaire}

\theoremstyle{remark}
\newtheorem{remarque}[definition]{Remarque}
\newtheorem{limitation}[definition]{Limitation}

% ----------------------------
% Notations
% ----------------------------

\newcommand{\D}{\mathcal{D}}
\newcommand{\B}{\mathcal{B}}
\newcommand{\G}{\mathcal{G}}
\newcommand{\C}{\mathcal{C}}
\newcommand{\E}{\mathcal{E}}
\newcommand{\T}{\mathcal{T}}
\newcommand{\A}{\mathcal{A}}
\newcommand{\R}{\mathbb{R}}
\newcommand{\N}{\mathbb{N}}
\newcommand{\X}{\mathcal{X}}
\newcommand{\F}{\mathcal{F}}
\newcommand{\I}{\mathcal{I}}
\newcommand{\Path}{\operatorname{Path}}
\newcommand{\Occ}{\operatorname{Occ}}
\newcommand{\supp}{\operatorname{supp}}
\newcommand{\id}{\operatorname{id}}
\newcommand{\Compat}{\operatorname{Compat}}
\newcommand{\Align}{\operatorname{Align}}
\newcommand{\Nat}{\operatorname{Nat}}
\newcommand{\Anti}{\operatorname{Anti}}
\newcommand{\Obs}{\operatorname{Obs}}
\newcommand{\FP}{\operatorname{FP}}
\newcommand{\Emb}{\operatorname{Emb}}
\newcommand{\Next}{\operatorname{Next}}
\newcommand{\Hist}{\operatorname{Hist}}
\newcommand{\End}{\operatorname{End}}

\title{
  \textbf{Théorie du maillage contextuel}\\
  \textbf{et de la navigation causale}\\[0.6em]
  \large Formalisation des brains, de la fusion des tokens,
  des écosystèmes relationnels, de l'alignement et de la recherche
  contextuelle de réponse
}
\author{}
\date{}

\begin{document}

\maketitle
\tableofcontents

% ============================================================
\chapter{Objet, périmètre et principes fondamentaux}
% ============================================================

\section{Objet de la théorie}

La présente théorie décrit un système de traitement du langage
construit directement à partir d'un dataset de couples
prompt/réponse, sans exiger qu'un modèle neuronal de type LLM
soit préalablement entraîné pour construire sa structure
relationnelle.

Le système suit quatre opérations principales :

\[
\boxed{
\text{Dataset}
\longrightarrow
\text{Brains}
\longrightarrow
\text{Maillage global}
\longrightarrow
\text{Navigation contextuelle}
}
\]

Le dataset est d'abord décomposé en exemples individuels.
Chaque exemple donne lieu à un brain qui conserve sa trajectoire
et sa provenance. Les identités communes des tokens sont ensuite
fusionnées dans un maillage global, sans supprimer les relations
qui proviennent des trajectoires d'origine.

Lorsqu'un nouveau prompt est présenté, un maillage contextuel
distinct est construit. Il sert à rechercher, dans le maillage
global, des trajectoires qui correspondent au prompt et qui
permettent de construire une réponse.

La cohérence recherchée repose conjointement sur :

\begin{enumerate}[label=\arabic*.]
  \item l'alignement entre le prompt et une structure du maillage
        global ;
  \item la possibilité d'intermédiaires entre les éléments alignés ;
  \item l'orientation de la trajectoire de recherche ;
  \item les directions naturelles conservées dans le maillage ;
  \item les relations anti-causales du maillage ;
  \item la compatibilité de la réponse en construction avec le
        prompt complet et avec les éléments déjà produits.
\end{enumerate}

La théorie distingue ces mécanismes de la vérité factuelle :
une structure cohérente peut contenir des informations erronées
si les données qui l'ont construite sont erronées.

\section{Principes non négociables}

\begin{postulat}[Autonomie initiale des brains]
Chaque exemple du dataset est d'abord représenté par un brain
propre. Les trajectoires provenant d'exemples distincts ne sont
pas confondues lors de cette première construction.
\end{postulat}

\begin{postulat}[Fusion des identités communes]
Le maillage global possède une identité canonique pour chaque
token reconnu comme commun selon la règle d'identité déclarée.
La fusion ne doit pas supprimer les relations ni la provenance
des occurrences qui ont été fusionnées.
\end{postulat}

\begin{postulat}[Séparation des deux maillages]
Le maillage contextuel d'un prompt n'est jamais incorporé au
maillage global du seul fait qu'il est construit. Il constitue
un motif de recherche indépendant qui interroge le maillage
global.
\end{postulat}

\begin{postulat}[Double condition de cohérence structurelle]
Une réponse recherchée doit satisfaire un alignement contextuel
et une organisation orientée de la trajectoire. La sélection
doit également tenir compte des relations naturelles et
anti-causales déjà inscrites dans le maillage global.
\end{postulat}

\begin{remarque}
Ces postulats définissent l'architecture théorique. Ils ne
constituent pas, à eux seuls, une preuve expérimentale que le
système atteint une qualité de réponse déterminée.
\end{remarque}

\section{Niveaux à ne pas confondre}

La théorie distingue les objets suivants :

\begin{itemize}
  \item la donnée brute : le texte fourni par le dataset ;
  \item l'occurrence : une apparition déterminée d'un token dans
        un exemple précis ;
  \item l'identité canonique : le token partagé auquel plusieurs
        occurrences sont rattachées ;
  \item le brain : la structure relationnelle d'un exemple ;
  \item le maillage global : le réseau obtenu par fusion des
        identités communes ;
  \item le maillage contextuel : la structure construite à partir
        du prompt courant ;
  \item l'alignement : la correspondance entre le motif contextuel
        et une structure du maillage global ;
  \item la réponse candidate : une trajectoire proposée par la
        recherche ;
  \item la réponse émise : la réalisation textuelle retenue par
        la politique de sortie.
\end{itemize}

Une relation entre tokens ne constitue pas automatiquement une
relation sémantique certifiée. Une direction appelée
« causale » dans le modèle ne constitue pas automatiquement une
causalité physique ou une preuve de vérité.

% ============================================================
\chapter{Dataset, tokenisation et construction des brains}
% ============================================================

\section{Dataset}

\begin{definition}[Dataset de couples prompt/réponse]
Un dataset est une famille finie ou dénombrable
\[
\D=\{d_1,d_2,\ldots,d_n\},
\]
où chaque exemple est de la forme
\[
d_i=(P_i,R_i,m_i).
\]
Ici :
\begin{itemize}
  \item \(P_i\) est le prompt ;
  \item \(R_i\) est la réponse associée ;
  \item \(m_i\) désigne les métadonnées disponibles, notamment
        l'identifiant de l'exemple, la source et les annotations
        éventuelles.
\end{itemize}
\end{definition}

La réponse \(R_i\) est conservée comme réponse observée dans le
dataset. Sa présence ne prouve pas, à elle seule, qu'elle est
vraie ou optimale.

\section{Tokenisation et identité}

\begin{definition}[Tokenisation]
Une fonction de tokenisation est une application
\[
\operatorname{Tok}:\text{Textes}\longrightarrow\Sigma^*,
\]
où \(\Sigma\) est l'ensemble des identités de tokens du système
et \(\Sigma^*\) l'ensemble des séquences finies de ces tokens.

Pour chaque exemple \(d_i\), on écrit
\[
\operatorname{Tok}(P_i)=(p_{i,1},\ldots,p_{i,k_i})
\]
et
\[
\operatorname{Tok}(R_i)=(r_{i,1},\ldots,r_{i,\ell_i}).
\]
\end{definition}

La règle d'identité des tokens doit être fixée avant la fusion.
Elle peut reposer sur une tokenisation donnée et une normalisation
explicitement déclarée.

Par défaut, deux tokens sont communs lorsqu'ils possèdent la
même identité lexicale canonique dans cette tokenisation.

Il ne faut pas confondre :
\[
\text{même forme de token}
\qquad\text{et}\qquad
\text{même sens dans tous les contextes}.
\]

La fusion porte sur l'identité déclarée, pas sur une hypothèse
implicite d'identité sémantique universelle.

\section{Occurrences et identité canonique}

\begin{definition}[Occurrence de token]
Une occurrence est un objet identifié par un quadruplet
\[
o=(i,s,j,x),
\]
où \(i\) identifie l'exemple, \(s\in\{P,R\}\) désigne le prompt
ou la réponse, \(j\) est la position dans la séquence et \(x\)
est l'identité canonique du token.

Deux occurrences peuvent donc être distinctes tout en partageant
la même identité canonique.
\end{definition}

Par exemple, les occurrences de \texttt{bien} dans
\[
\texttt{je vais bien},
\qquad
\texttt{il va très bien aujourd'hui},
\qquad
\texttt{c'est bien}
\]
sont trois occurrences distinctes, mais elles peuvent être
rattachées au même nœud canonique \texttt{bien}.

Cette distinction est fondamentale :
la fusion rend l'identité commune, mais ne doit pas effacer
l'histoire des occurrences.

\section{Trajectoire d'un exemple}

\begin{definition}[Trajectoire ordonnée d'un exemple]
Pour un exemple \(d_i\), la trajectoire textuelle complète est
la séquence ordonnée
\[
\tau_i=
(p_{i,1},\ldots,p_{i,k_i},
r_{i,1},\ldots,r_{i,\ell_i}).
\]
Elle conserve séparément la frontière entre prompt et réponse.
\end{definition}

On représente les relations de succession observées dans le
prompt par
\[
p_{i,j}\longrightarrow p_{i,j+1},
\qquad 1\leq j<k_i,
\]
et celles observées dans la réponse par
\[
r_{i,j}\longrightarrow r_{i,j+1},
\qquad 1\leq j<\ell_i.
\]

La frontière prompt/réponse est représentée par une relation
typée :
\[
p_{i,k_i}
\overset{\mathrm{début\_réponse}}{\longrightarrow}
r_{i,1}.
\]

Cette dernière relation signifie que la réponse suit le prompt
dans l'exemple observé. Elle ne signifie pas que le dernier token
du prompt est la cause physique de la réponse.

\section{Définition du brain}

\begin{definition}[Brain d'un exemple]
Le brain de l'exemple \(d_i\) est un multigraphe orienté, typé et
traçable
\[
B_i=(V_i,E_i,\lambda_i,\pi_i).
\]
Ses composantes sont :

\begin{itemize}
  \item \(V_i\) : les nœuds représentant les occurrences de tokens
        de l'exemple ;
  \item \(E_i\) : les relations orientées entre ces occurrences ;
  \item \(\lambda_i\) : l'étiquette de chaque relation, par
        exemple succession dans le prompt, succession dans la
        réponse ou frontière prompt/réponse ;
  \item \(\pi_i\) : la provenance des nœuds et des relations,
        notamment l'identifiant de l'exemple et les positions.
\end{itemize}
\end{definition}

Le brain n'est pas un simple ensemble de tokens. Il conserve
l'ordre observé, les frontières entre prompt et réponse et
l'origine de chaque relation.

\begin{proposition}[Conservation de la trajectoire dans un brain]
Si le brain conserve les occurrences, leur ordre et les relations
de succession, la trajectoire textuelle de l'exemple peut être
reconstruite à partir de sa structure ordonnée.
\end{proposition}

\begin{proof}
Chaque occurrence conserve sa position dans l'exemple et chaque
relation de succession relie deux positions consécutives.
La succession des occurrences est donc reconstructible à partir
des relations et des positions conservées.
\end{proof}

\begin{remarque}
Cette proposition concerne la reconstruction de la trajectoire
observée. Elle ne garantit pas que la trajectoire contient toutes
les informations sémantiques ou toutes les causes réelles de
l'énoncé.
\end{remarque}

% ============================================================
\chapter{Fusion des brains et construction du maillage global}
% ============================================================

\section{Principe de fusion}

Au départ, les brains sont construits séparément :
\[
B_1,B_2,\ldots,B_n.
\]

La fusion intervient ensuite sur les identités communes.
Elle ne consiste pas à supprimer arbitrairement les trajectoires
pour ne conserver qu'une liste de mots.

Le principe est :
\[
\boxed{
\text{identité canonique commune}
+
\text{relations conservées}
+
\text{provenance conservée}
}
\]

\section{Définition du maillage global}

\begin{definition}[Maillage global]
Le maillage global est un multigraphe orienté, typé et traçable
\[
G=(V_G,E_G,\lambda_G,\pi_G).
\]

Il est construit à partir de l'ensemble des brains. Les occurrences
qui partagent une même identité canonique sont rattachées à un
nœud canonique commun, tandis que les relations issues des brains
sont conservées avec leur type et leur provenance.
\end{definition}

Pour une identité canonique \(x\), on note
\[
v_x\in V_G
\]
le nœud correspondant.

L'ensemble des occurrences qui lui sont rattachées est
\[
\Occ(x)=
\{o:\text{l'identité canonique de }o\text{ est }x\}.
\]

Les relations issues des occurrences sont projetées sur les
identités canoniques, sans supprimer leurs traces d'origine.

\section{Fusion et préservation des relations}

Soit
\[
q: \bigcup_{i=1}^{n}V_i\longrightarrow V_G
\]
l'application qui rattache chaque occurrence à son identité
canonique.

Si une relation orientée de type \(\ell\) relie une occurrence
\(u\) à une occurrence \(v\), la relation correspondante dans
le maillage global est
\[
q(u)\overset{\ell}{\longrightarrow}q(v),
\]
avec conservation de la provenance de cette relation.

Plusieurs relations peuvent aboutir à la même paire de nœuds
canoniques. Elles ne doivent pas être confondues si elles diffèrent
par leur type, leur contexte ou leur provenance.

\begin{proposition}[La fusion conserve les relations tracées]
Si chaque relation fusionnée conserve son type et sa provenance,
alors la fusion des identités communes ne supprime pas
l'information permettant d'identifier les relations observées
dans les brains d'origine.
\end{proposition}

\begin{proof}
Toute relation d'un brain est projetée vers une relation du
maillage global et conserve une référence à son brain d'origine.
La relation peut ainsi être retrouvée depuis le maillage global
par sa provenance. La fusion peut donc partager les identités
sans effacer les trajectoires sources.
\end{proof}

\begin{remarque}
La proposition dépend de la conservation effective des références
de provenance. Si les relations sont fusionnées en une simple
arête non typée, sans historique ni multiplicité, cette propriété
n'est plus garantie.
\end{remarque}

\section{La redondance utile}

La redondance d'identité est réduite :
plusieurs occurrences du même token partagent un nœud canonique.

En revanche, la diversité des relations reste présente.

Ainsi, le nœud \texttt{bien} n'est pas une occurrence unique
qui aurait perdu son histoire. Il devient un point d'accès commun
à plusieurs relations et à plusieurs trajectoires observées.

On peut schématiser :
\[
\begin{aligned}
B_1 &: \texttt{je}\to\texttt{vais}\to\texttt{bien},\\
B_2 &: \texttt{il}\to\texttt{va}\to\texttt{très}
       \to\texttt{bien}\to\texttt{aujourd'hui},\\
B_3 &: \texttt{c'est}\to\texttt{bien}.
\end{aligned}
\]

Après fusion, les occurrences de \texttt{bien} partagent une
identité canonique. Les relations provenant des trois brains
restent néanmoins distinguables par leur direction, leur type
et leur provenance.

\section{Directions naturelles}

\begin{definition}[Direction naturelle observée]
Une direction naturelle observée est une relation orientée
\[
x\overset{\mathrm{Nat}}{\longrightarrow}y
\]
qui provient d'une succession observée dans une trajectoire,
ou d'une règle explicite du modèle déclarée comme naturelle.
\end{definition}

Dans une réponse telle que
\[
\texttt{je vais bien merci},
\]
la succession observée donne notamment
\[
\texttt{bien}
\overset{\mathrm{Nat}}{\longrightarrow}
\texttt{merci}.
\]

Dans
\[
\texttt{c'est bien de le dire},
\]
on observe notamment
\[
\texttt{bien}
\overset{\mathrm{Nat}}{\longrightarrow}
\texttt{de}.
\]

Dans
\[
\texttt{oui bien que},
\]
on observe notamment
\[
\texttt{bien}
\overset{\mathrm{Nat}}{\longrightarrow}
\texttt{que}.
\]

Le nœud \texttt{bien} dispose donc de plusieurs directions
naturelles issues de trajectoires distinctes.

\section{L'écosystème d'un token}

\begin{definition}[Écosystème relationnel]
L'écosystème relationnel d'un token canonique \(x\), noté
\[
\E(x),
\]
est l'ensemble des relations du maillage global incidentes à
\(x\), conservées avec leur orientation, leur type et leur
provenance.
\end{definition}

L'écosystème ne se réduit pas aux voisins immédiatement
suivants dans une seule phrase. Il rassemble les relations
issues de l'ensemble des brains dans lesquels l'identité
canonique apparaît.

On distingue :
\[
\E_{\mathrm{Nat}}(x)
\]
pour les relations naturelles, et
\[
\E_{\mathrm{Anti}}(x)
\]
pour les relations anti-causales.

Ces ensembles sont des sous-ensembles typés de l'écosystème
global :
\[
\E(x)=
\E_{\mathrm{Nat}}(x)
\cup
\E_{\mathrm{Anti}}(x)
\cup
\E_{\mathrm{autres}}(x).
\]

\begin{remarque}
L'écosystème est plus riche qu'une liste de voisins : il contient
les types de relations et leurs traces. Deux tokens peuvent avoir
les mêmes voisins apparents tout en possédant des écosystèmes
différents si les directions, les contextes ou les provenances
diffèrent.
\end{remarque}

\section{Directions naturelles et multiplicité}

Le nombre d'occurrences d'un token n'est pas l'unique information
conservée par le maillage. La diversité des directions observées
est elle-même une partie de l'information relationnelle.

Pour chaque paire de tokens \((x,y)\), on peut conserver :
\[
M_{\mathrm{Nat}}(x,y)
=
\{\text{traces des relations naturelles }x\to y\}.
\]

La multiplicité observable de la relation peut être calculée par
\[
m_{\mathrm{Nat}}(x,y)
=
|M_{\mathrm{Nat}}(x,y)|
\]
lorsque cet ensemble est fini.

Cette multiplicité est une mesure de fréquence ou de répétition
des relations enregistrées. Elle ne doit pas être confondue avec
la totalité de l'information contextuelle.

\begin{proposition}[La fusion n'oblige pas à perdre la multiplicité]
Le partage d'un nœud canonique par plusieurs occurrences est
compatible avec la conservation du nombre de relations observées,
dès lors que les relations et leurs provenances ne sont pas
supprimées.
\end{proposition}

\begin{proof}
La multiplicité est calculée sur l'ensemble des traces de
relations conservées. La fusion des identités ne supprime pas
ces traces lorsque la règle de fusion les préserve. Le nombre
de relations observées reste donc calculable.
\end{proof}

% ============================================================
\chapter{Causalité naturelle et anti-causalité}
% ============================================================

\section{Statut de la causalité dans cette théorie}

Le mot « causalité » désigne ici d'abord une structure interne
du modèle : une orientation et un type de relation utilisés pour
organiser les trajectoires et contraindre la recherche.

Il ne faut pas automatiquement interpréter une relation
\[
x\to y
\]
comme la preuve que \(x\) est une cause physique ou sémantique
de \(y\).

Une relation observée dans un texte peut représenter un ordre,
une continuation, une association contextuelle ou une annotation.
Une affirmation de causalité réelle demanderait des éléments
supplémentaires qui ne sont pas fournis par la seule succession
des tokens.

\section{Causalité naturelle}

\begin{definition}[Relation causale naturelle du modèle]
Une relation
\[
x\overset{\mathrm{Nat}}{\longrightarrow}y
\]
est une relation que le modèle classe dans sa catégorie
naturelle. Dans la construction initiale, elle est notamment
produite par la succession effective \(x,y\) dans une trajectoire
observée.
\end{definition}

La catégorie naturelle sert à préserver les continuations
observées et à les rendre accessibles dans le maillage global.

Exemples :
\[
\begin{aligned}
\texttt{bien}&\overset{\mathrm{Nat}}{\longrightarrow}\texttt{merci},\\
\texttt{bien}&\overset{\mathrm{Nat}}{\longrightarrow}\texttt{de},\\
\texttt{bien}&\overset{\mathrm{Nat}}{\longrightarrow}\texttt{que}.
\end{aligned}
\]

Ces trois relations proviennent de trajectoires textuelles
distinctes, mais elles partagent le même nœud canonique source.

\section{Anti-causalité}

\begin{definition}[Relation anti-causale du modèle]
Une relation
\[
x\overset{\mathrm{Anti}}{\longrightarrow}y
\]
est une relation explicitement classée comme anti-causale dans
le maillage. Elle conserve sa propre orientation et sa propre
étiquette. Elle n'est pas automatiquement identifiée à la
relation inverse \(y\to x\).
\end{definition}

Cette définition respecte le sens spécifique que la catégorie
anti-causale possède dans le modèle. Elle ne suppose pas que
l'anti-causalité soit simplement une inversion de l'ordre des
tokens.

Dans l'exemple de travail, on conserve les relations déclarées :
\[
\begin{aligned}
\texttt{bien}&\overset{\mathrm{Anti}}{\longrightarrow}\texttt{c'est},\\
\texttt{bien}&\overset{\mathrm{Anti}}{\longrightarrow}\texttt{oui},\\
\texttt{bien}&\overset{\mathrm{Anti}}{\longrightarrow}\texttt{vais}.
\end{aligned}
\]

Ces relations sont distinctes des relations naturelles
\(\texttt{bien}\to\texttt{merci}\),
\(\texttt{bien}\to\texttt{de}\) et
\(\texttt{bien}\to\texttt{que}\).

\begin{remarque}
Une relation anti-causale ne peut pas être déduite uniquement
de la présence d'un token avant un autre. Sa règle de
classification doit être définie et traçable dans le système.
Les exemples ci-dessus sont des étiquettes fournies par le
modèle, et non une loi universelle de la langue française.
\end{remarque}

\section{Fonction des relations anti-causales dans la recherche}

L'anti-causalité doit contribuer au mécanisme de sélection ;
elle ne peut pas être une simple étiquette décorative.

On distingue deux usages compatibles, dont l'application dépend
de la règle de recherche déclarée.

\begin{definition}[Exclusion anti-causale]
Une relation anti-causale joue un rôle d'exclusion lorsqu'une
règle déclarée établit qu'elle est incompatible avec la
continuation naturelle recherchée dans le contexte courant.
Dans ce cas, elle empêche de traiter cette relation comme une
continuation naturelle admissible.
\end{definition}

\begin{definition}[Réorientation anti-causale]
Une relation anti-causale joue un rôle de réorientation lorsqu'elle
signale une direction de contraste ou une région relationnelle
alternative à examiner. Elle ne devient pas pour autant une
continuation naturelle.
\end{definition}

\begin{postulat}[Usage opératoire de l'anti-causalité]
Lors de la recherche, une relation anti-causale doit influencer
soit l'admissibilité d'une continuation, soit l'orientation de
l'exploration, soit les deux lorsque la politique déclarée le
permet. Elle ne doit jamais être silencieusement convertie en
relation naturelle.
\end{postulat}

\begin{remarque}
La catégorie anti-causale ne signifie pas que chaque relation
anti-causale doit être rejetée dans toutes les situations.
Une relation peut être pertinente comme contraste ou comme
indication d'une autre région du maillage. La règle exacte
doit être déterminée par le protocole d'usage de la catégorie.
\end{remarque}

\section{Pourquoi les trois niveaux de contraintes coopèrent}

Le mécanisme global articule :
\[
\boxed{
\text{alignement du prompt}
+
\text{orientation de recherche}
+
\text{écosystème naturel/anti-causal}
}
\]

L'alignement repère les structures qui correspondent au prompt.
L'orientation détermine le sens du parcours de recherche.
L'écosystème distingue les continuations naturelles des relations
qui possèdent un autre statut.

Ces contraintes peuvent éliminer de nombreux parcours
incompatibles. Elles ne garantissent cependant pas que toute
réponse retenue soit factuellement vraie : cette propriété
dépend notamment de la qualité des données et des relations
enregistrées.

% ============================================================
\chapter{Construction du maillage contextuel du prompt}
% ============================================================

\section{Séparation structurelle}

\begin{definition}[Maillage contextuel]
Pour un prompt \(P\), le maillage contextuel est une structure
\[
C_P=(V_P,E_P,\lambda_P),
\]
construite à partir des tokens du prompt et de leurs relations
contextuelles. Il sert de motif de recherche dans le maillage
global.
\end{definition}

Le maillage contextuel n'est pas une partie du maillage global.
Il n'ajoute aucun token, aucune relation et aucune provenance
au maillage global par sa seule construction.

On dispose donc de deux objets distincts :
\[
G=(V_G,E_G,\lambda_G,\pi_G)
\]
et
\[
C_P=(V_P,E_P,\lambda_P).
\]

La recherche utilise \(C_P\) pour interroger \(G\), mais ne
fusionne pas leurs structures :
\[
C_P\not\hookrightarrow G
\]
au sens d'une insertion structurelle automatique.

\section{Tokens d'ancrage et variables contextuelles}

Soit le prompt tokenisé
\[
P=(x_1,x_2,\ldots,x_n).
\]

Les \(x_i\) sont les tokens d'ancrage du prompt. On intercale
entre eux des variables contextuelles
\[
c_0,c_1,\ldots,c_n,c_x.
\]

Les variables \(c_i\) ne désignent pas un token lexical fixé.
Elles désignent des emplacements de correspondance capables
d'accepter des structures du maillage global.

\begin{definition}[Valeur admissible d'une variable contextuelle]
Une valeur admissible d'une variable \(c_i\) peut être :
\begin{enumerate}[label=\alph*)]
  \item un token canonique ;
  \item une séquence finie de tokens ;
  \item une trajectoire orientée du maillage global ;
  \item un sous-graphe ou un ensemble de tokens muni de relations,
        lorsque la règle de correspondance autorise cette forme.
\end{enumerate}
\end{definition}

La liberté de \(c_i\) n'est pas une permission d'ignorer les
contraintes. La structure affectée à \(c_i\) doit appartenir au
maillage global et respecter les règles de compatibilité
déclarées.

\section{Colonne vertébrale ordonnée du prompt}

La structure de base du motif est :
\[
c_0\to x_1\to c_1\to x_2\to c_2
\to\cdots\to x_n\to c_x.
\]

Cette écriture exprime une organisation ordonnée des ancres et
des emplacements contextuels. Elle ne signifie pas que chaque
\(c_i\) correspond à un seul token ou à une seule arête.

Par exemple, pour
\[
P=\texttt{je vais bien},
\]
on écrit :
\[
c_0\to\texttt{je}\to c_1
\to\texttt{vais}\to c_2
\to\texttt{bien}\to c_x.
\]

Si le maillage global contient une trajectoire compatible dans
laquelle plusieurs tokens ou une trajectoire entière se trouvent
entre \texttt{je} et \texttt{vais}, la variable \(c_1\) peut
correspondre à cette structure intermédiaire. La même règle
s'applique aux autres variables contextuelles.

\section{Multidirectionnalité contextuelle}

La colonne vertébrale ne suffit pas à représenter l'idée du
maillage contextuel. Le prompt doit également être traité comme
une configuration relationnelle où les tokens peuvent contribuer
conjointement à l'orientation de la recherche.

On ajoute donc des relations contextuelles entre les ancres,
ainsi que des relations entre les ancres et les emplacements
contextuels.

Pour le prompt \texttt{je vais bien}, le motif peut contenir
notamment :
\[
\begin{aligned}
c_0&\to\texttt{je},&
c_0&\to\texttt{vais},&
c_0&\to\texttt{bien},&
c_0&\to c_x,\\
\texttt{je}&\to\texttt{vais},&
\texttt{je}&\to\texttt{bien},&
\texttt{je}&\to c_x,\\
\texttt{vais}&\to\texttt{je},&
\texttt{vais}&\to\texttt{bien},&
\texttt{vais}&\to c_x,\\
\texttt{bien}&\to\texttt{je},&
\texttt{bien}&\to\texttt{vais},&
\texttt{bien}&\to c_x.
\end{aligned}
\]

Les relations précises du motif sont fixées par la règle de
construction contextuelle. Le motif peut conserver des relations
réciproques entre tokens même lorsque sa colonne vertébrale
principale est ordonnée.

\begin{remarque}
Une relation contextuelle réciproque n'est pas automatiquement
une relation naturelle du maillage global. Les relations du
motif décrivent comment le prompt est utilisé comme contexte ;
les relations du maillage global décrivent la structure
relationnelle enregistrée dans le dataset.
\end{remarque}

\section{La direction \(c_0\to c_x\)}

\begin{definition}[Orientation de recherche]
La relation
\[
c_0\longrightarrow c_x
\]
désigne l'orientation globale du parcours de recherche associé
au motif contextuel. Elle organise la progression depuis
l'entrée contextuelle vers la sortie recherchée.
\end{definition}

Cette orientation ne doit pas être confondue avec :
\begin{itemize}
  \item l'ordre de lecture seul ;
  \item une relation réciproque de contextualisation ;
  \item une causalité physique démontrée.
\end{itemize}

Dans le modèle, \(c_0\to c_x\) impose que la recherche soit
orientée. Les variables intermédiaires permettent de conserver
la souplesse de l'alignement tout en gardant cette orientation.

% ============================================================
\chapter{Alignement contextuel dans le maillage global}
% ============================================================

\section{Définition de l'alignement}

\begin{definition}[Alignement contextuel]
Un alignement du prompt \(P\) dans le maillage global \(G\) est
une correspondance entre les ancres du maillage contextuel et
des occurrences ou identités du maillage global, accompagnée
d'une affectation admissible aux variables \(c_i\), qui respecte
les contraintes structurelles déclarées du motif.
\end{definition}

On note l'ensemble des alignements admissibles :
\[
\Align(P,G).
\]

Un élément
\[
\alpha\in\Align(P,G)
\]
est un certificat d'alignement. Il précise au minimum :
\begin{itemize}
  \item où chaque ancre du prompt est retrouvée ;
  \item quelle structure est affectée à chaque variable \(c_i\) ;
  \item quels chemins ou relations relient les ancres ;
  \item quelle orientation est utilisée ;
  \item quelles relations et provenances du maillage global
        soutiennent la correspondance.
\end{itemize}

\section{Les intermédiaires sont autorisés}

Pour un prompt abstrait
\[
P=(a,b,c),
\]
un alignement ne doit pas exiger que \(a\), \(b\) et \(c\)
soient trois voisins directs.

Il peut avoir la forme :
\[
c_0\to a\to
\underbrace{u_1\to\cdots\to u_k}_{c_1}
\to b\to
\underbrace{v_1\to\cdots\to v_m}_{c_2}
\to c\to c_x.
\]

Les tokens \(u_j\) et \(v_j\) sont des éléments du maillage
global qui servent d'intermédiaires. Ils ne sont pas supposés
faire partie du prompt original.

Les variables contextuelles peuvent aussi correspondre à une
structure relationnelle plus large lorsque la règle de
correspondance le permet.

\begin{definition}[Alignement ordonné avec intermédiaires]
Un alignement est ordonné avec intermédiaires lorsque les ancres
\(a,b,c\) sont retrouvées dans l'ordre exigé par le motif, tout
en autorisant entre elles des chemins ou sous-structures
admissibles du maillage global.
\end{definition}

\section{Correspondance et identité lexicale}

La correspondance entre un token du prompt et un token canonique
du maillage global repose sur la règle d'identité déclarée.

Si \(x\) est un token du prompt, on recherche un nœud global
\(v_x\) portant la même identité canonique, ou une correspondance
explicitement autorisée par le protocole.

Cette correspondance ne fusionne pas le nœud du prompt avec le
nœud global. Elle est une relation de recherche :
\[
\mu_P:V_P\longrightarrow V_G
\]
lorsqu'une telle correspondance est définie.

La fonction \(\mu_P\) est une application de correspondance,
pas une opération de fusion structurelle.

\section{Pourquoi la présence des tokens ne suffit pas}

\begin{proposition}[La présence isolée des ancres n'implique pas
un alignement admissible]
La présence de chacun des tokens d'un prompt dans le maillage
global ne suffit pas à établir un alignement contextuel. Il faut
également satisfaire les contraintes de relation, d'ordre et
d'orientation du motif.
\end{proposition}

\begin{proof}
L'appartenance individuelle des ancres à \(V_G\) n'établit
aucune relation entre leurs occurrences. Or la définition de
l'alignement exige une affectation des variables contextuelles
et des relations qui satisfont le motif. La présence des tokens
est donc une condition éventuellement nécessaire, mais non
suffisante.
\end{proof}

\section{La condition d'orientation}

Un alignement doit respecter l'orientation de recherche
\(c_0\to c_x\). Les correspondances ne sont pas acceptées
uniquement parce qu'elles contiennent les bons tokens.

\begin{definition}[Alignement orienté]
Un alignement est orienté lorsque les structures affectées aux
variables contextuelles et les relations utilisées permettent
de parcourir le motif dans le sens \(c_0\to c_x\), conformément
aux contraintes déclarées.
\end{definition}

L'orientation est une condition sur le parcours de recherche.
Elle n'établit pas, par elle-même, que la réponse est vraie.

% ============================================================
\chapter{Recherche de réponse et activation du maillage}
% ============================================================

\section{Du motif à la région pertinente}

Le maillage contextuel ne sert pas uniquement à attribuer un
score aux tokens du maillage global. Il définit une configuration
à rechercher.

Pour un prompt \(P\), l'ensemble
\[
\Align(P,G)
\]
réunit les alignements admissibles.

Chaque alignement active une partie du maillage global :
les ancres correspondantes, les chemins intermédiaires et les
relations qui soutiennent l'alignement.

\begin{definition}[Région activée]
La région activée par le prompt est l'union des nœuds et relations
qui participent à au moins un alignement admissible :
\[
\A(P,G)=
\bigcup_{\alpha\in\Align(P,G)}
\operatorname{Support}(\alpha).
\]
Si aucun alignement n'existe, on pose
\[
\A(P,G)=\varnothing.
\]
\end{definition}

Cette définition exprime le fait que le contexte peut activer
plusieurs chemins compatibles à la fois. Elle ne réduit pas
automatiquement le contexte à un seul chemin gagnant.

\section{Diversité compatible avec le contexte}

La diversité des directions est prise en compte à travers
l'ensemble des alignements admissibles et leurs supports.
Une direction ne devient pas pertinente uniquement parce
qu'elle est fréquente ; elle doit appartenir à une structure
compatible avec le prompt.

Pour un nœud \(x\), on peut définir l'ensemble des directions
naturelles activées par le prompt :
\[
\Next_{\mathrm{Nat}}(x\mid P,G)=
\left\{
y:
x\overset{\mathrm{Nat}}{\longrightarrow}y
\text{ et la relation appartient à une continuation compatible}
\right\}.
\]

De même, les relations anti-causales activées sont :
\[
\Next_{\mathrm{Anti}}(x\mid P,G)=
\left\{
y:
x\overset{\mathrm{Anti}}{\longrightarrow}y
\text{ et la relation est pertinente pour la recherche courante}
\right\}.
\]

Ces deux ensembles ne sont pas interchangeables. Le premier
contient des directions naturelles admissibles ; le second
contient des relations anti-causales qui peuvent contribuer
à l'exclusion ou à la réorientation selon le protocole.

\section{Rôle du prompt complet}

Le choix d'une continuation ne dépend pas seulement du dernier
token atteint. Il dépend du prompt complet et de la structure
déjà parcourue.

On note \(H_t\) l'historique de la réponse produite jusqu'à
l'étape \(t\), et \(S_t\) l'état de recherche, qui contient
notamment :
\begin{itemize}
  \item le motif contextuel courant ;
  \item les ancres déjà alignées ;
  \item les chemins déjà parcourus ;
  \item les relations naturelles et anti-causales examinées ;
  \item les éléments de réponse déjà produits ;
  \item les certificats de provenance des continuations candidates.
\end{itemize}

La sélection suivante s'écrit abstraitement :
\[
S_{t+1}=
\mathcal{U}(S_t,G,C_P,H_t),
\]
où \(\mathcal{U}\) est la règle de mise à jour déclarée.

La génération est donc une navigation dynamique : chaque
nouvel élément produit modifie l'état de recherche et peut
modifier la région activée.

\section{Recherche d'une réponse dans les trajectoires du dataset}

Les brains conservent la frontière entre prompt et réponse.
Le maillage global conserve les relations qui en proviennent.
Lorsqu'un alignement du prompt est trouvé, le système peut
rechercher les continuations de réponse soutenues par les
trajectoires sources correspondantes.

On distingue deux voies compatibles :

\begin{enumerate}[label=\arabic*.]
  \item \textbf{Voie par exemple aligné :} l'alignement correspond
        à une région d'un brain source. Le système examine la
        réponse attachée à cet exemple et les relations qui
        permettent de la retrouver.
  \item \textbf{Voie par navigation relationnelle :} plusieurs
        relations du maillage global sont combinées pour trouver
        une continuation qui respecte l'alignement, l'orientation
        et l'écosystème naturel/anti-causal.
\end{enumerate}

La seconde voie autorise une réponse qui ne reproduit pas
nécessairement une seule réponse mémorisée, à condition que
la continuation construite reste soutenue par les relations
et les trajectoires disponibles.

\section{Condition structurelle de recevabilité d'une réponse}

\begin{definition}[Réponse structurellement recevable]
Une réponse candidate \(R\) est structurellement recevable pour
un prompt \(P\) relativement au maillage global \(G\) si :

\begin{enumerate}[label=(\roman*)]
  \item il existe au moins un alignement admissible
        \(\alpha\in\Align(P,G)\) ;
  \item la trajectoire de recherche respecte l'orientation
        \(c_0\to c_x\) ;
  \item la continuation proposée est reliée à la région alignée
        par des relations et des trajectoires admissibles ;
  \item les relations naturelles et anti-causales pertinentes
        ont été prises en compte conformément à leur type ;
  \item la provenance des relations utilisées est conservée.
\end{enumerate}
\end{definition}

\begin{proposition}[Absence de soutien structurel]
Si aucune relation ni trajectoire admissible ne relie une réponse
candidate à un alignement du prompt, cette réponse ne satisfait
pas la définition de recevabilité structurelle.
\end{proposition}

\begin{proof}
La définition exige une continuation reliée à une région alignée
par des relations et des trajectoires admissibles. Si un tel
soutien n'existe pas, la condition (iii) est fausse. La réponse
ne peut donc pas être déclarée structurellement recevable.
\end{proof}

Cette proposition formalise le rejet d'une réponse arbitraire
sans correspondance ni chemin de soutien. Elle ne prétend pas
que toute réponse soutenue structurellement est vraie.

% ============================================================
\chapter{Boucles, répétitions et progression de la recherche}
% ============================================================

\section{Le problème des boucles}

Un maillage peut contenir des cycles. Par exemple :
\[
\texttt{bien}\to\texttt{très}\to\texttt{bien}.
\]

Si la recherche se limite à suivre le voisinage local du dernier
token, elle peut répéter indéfiniment ce cycle sans faire
progresser la réponse.

La solution ne consiste pas à interdire toute répétition de
token : une répétition peut être légitime dans une langue.
Il faut distinguer une répétition textuelle utile d'une boucle
de recherche qui ne fait aucun progrès.

\section{État de progression}

\begin{definition}[État de progression]
L'état de progression \(S_t\) comprend la position de recherche
dans le motif contextuel, l'historique \(H_t\), les alignements
retenus, les relations déjà examinées et les conditions de
terminaison.
\end{definition}

\begin{postulat}[Progression ou terminaison]
Chaque étape de recherche doit satisfaire au moins l'une des
conditions suivantes :
\begin{enumerate}[label=\alph*)]
  \item elle fait progresser l'alignement ou la continuation
        vers une partie non encore satisfaite du motif ;
  \item elle apporte une nouvelle relation ou un nouveau
        certificat de soutien pertinent ;
  \item elle explore explicitement une branche alternative ;
  \item elle termine la recherche selon une règle de terminaison
        déclarée.
\end{enumerate}
Une transition qui ne satisfait aucune de ces conditions ne
peut pas être répétée indéfiniment comme si elle apportait un
progrès.
\end{postulat}

Cette règle ne supprime pas les cycles du maillage. Elle évite
que la procédure de recherche reste bloquée dans un cycle sans
avancer dans la tâche.

\section{Terminaison}

La réponse peut être terminée lorsqu'une condition de sortie
déclarée est satisfaite, notamment :
\begin{itemize}
  \item rencontre d'une frontière de fin de réponse conservée
        dans les données ;
  \item atteinte d'une condition de complétude du motif de sortie ;
  \item atteinte d'une limite de génération fixée par le système ;
  \item absence de continuation admissible, auquel cas le système
        doit signaler l'échec ou l'insuffisance de soutien plutôt
        que fabriquer silencieusement une réponse.
\end{itemize}

La limite de génération est une contrainte d'implémentation.
Elle ne doit pas être présentée comme une propriété intrinsèque
du maillage.

% ============================================================
\chapter{Procédure opérationnelle complète}
% ============================================================

\section{Construction hors ligne}

La construction du maillage global s'effectue en plusieurs
étapes.

\begin{enumerate}[label=\arabic*.]
  \item Tokeniser chaque prompt et chaque réponse avec la même
        règle d'identité déclarée.
  \item Construire un brain par exemple, en conservant les
        occurrences, l'ordre, les frontières prompt/réponse et
        les relations de succession.
  \item Déterminer les identités canoniques communes selon la
        règle de normalisation fixée.
  \item Fusionner les identités communes dans le maillage global.
  \item Conserver toutes les relations utiles avec leur type,
        leur orientation et leur provenance.
  \item Distinguer les directions naturelles des relations
        anti-causales selon une règle d'étiquetage explicite.
  \item Vérifier que chaque relation du maillage global peut être
        reliée à une provenance ou à une règle déclarée.
\end{enumerate}

\section{Recherche en ligne}

Pour un nouveau prompt \(P\), la procédure est :

\begin{enumerate}[label=\arabic*.]
  \item Construire le maillage contextuel \(C_P\), sans l'insérer
        dans le maillage global.
  \item Identifier les ancres et les relations contextuelles
        du motif.
  \item Rechercher les correspondances des ancres dans le
        maillage global.
  \item Affecter les variables \(c_i\) à des tokens, chemins ou
        structures admissibles.
  \item Retenir les alignements qui respectent les contraintes
        de relation et l'orientation \(c_0\to c_x\).
  \item Activer les régions du maillage global soutenant ces
        alignements.
  \item Examiner les continuations issues des trajectoires
        alignées et des relations naturelles disponibles.
  \item Utiliser les relations anti-causales pour l'exclusion ou
        la réorientation, selon la règle déclarée.
  \item Mettre à jour l'état contextuel en tenant compte du prompt
        complet et de l'historique de réponse.
  \item Continuer tant qu'une progression admissible est possible.
  \item Terminer sur une condition explicite, ou signaler qu'aucune
        continuation suffisamment soutenue n'a été trouvée.
\end{enumerate}

\section{Pseudo-code de référence}

Le pseudo-code suivant décrit la logique, sans imposer un langage
de programmation ni une structure de données particulière.

\begin{verbatim}
CONSTRUCTION_HORS_LIGNE(dataset):
    brains = []

    pour chaque exemple (prompt, reponse, metadata):
        tokens_prompt  = TOKENISER(prompt)
        tokens_reponse = TOKENISER(reponse)

        brain = CONSTRUIRE_BRAIN(
            tokens_prompt,
            tokens_reponse,
            metadata
        )

        brains.ajouter(brain)

    G = FUSIONNER_IDENTITES(brains)
    CONSERVER_RELATIONS_TYPES_PROVENANCES(G)
    ETIQUETER_RELATIONS_NATURELLES_ET_ANTICAUSALES(G)

    retourner G


RECHERCHE(P, G):
    C = CONSTRUIRE_MAILLAGE_CONTEXTUEL(P)
    # C reste distinct de G.

    alignements = TROUVER_ALIGNEMENTS(C, G)

    si alignements est vide:
        retourner ECHEC_ALIGNEMENT

    region = ACTIVER_SUPPORTS(alignements)
    etat = INITIALISER_ETAT(P, C, region)

    tant que non TERMINAL(etat):
        candidats = TROUVER_CONTINUATIONS(etat, G)

        candidats = APPLIQUER_ORIENTATION_C0_VERS_CX(candidats)
        candidats = VERIFIER_COMPATIBILITE_NATURELLE(candidats)
        candidats = TRAITER_RELATIONS_ANTICAUSALES(candidats)

        si candidats est vide:
            retourner ECHEC_CONTINUATION

        etat = METTRE_A_JOUR_ETAT(
            etat,
            candidats,
            prompt_complet=P,
            historique=etat.historique
        )

        si non PROGRESSION(etat):
            retourner ECHEC_PROGRESSION

    retourner REALISER_REPONSE(etat)
\end{verbatim}

\begin{remarque}
Les opérations \texttt{TROUVER\_ALIGNEMENTS},
\texttt{VERIFIER\_COMPATIBILITE\_NATURELLE} et
\texttt{TRAITER\_RELATIONS\_ANTICAUSALES} sont des modules
dont les règles doivent être définies et testées. Le pseudo-code
ne prétend pas que leur simple nomination constitue déjà un
algorithme exécutable complet.
\end{remarque}

% ============================================================
\chapter{Exemple intégré : l'écosystème de \texttt{bien}}
% ============================================================

\section{Brains initiaux}

Considérons les exemples suivants :
\[
\begin{aligned}
P_1/R_1 &: \texttt{je vais bien normalement},\\
P_2/R_2 &: \texttt{il va très bien aujourd'hui},\\
P_3/R_3 &: \texttt{c'est bien de le dire},\\
P_4/R_4 &: \texttt{oui bien que},\\
P_5/R_5 &: \texttt{il faut beau ce soir}.
\end{aligned}
\]

Chaque ligne représente ici une trajectoire textuelle illustrative.
Dans un dataset réel, les frontières entre prompts et réponses
doivent être conservées séparément, conformément à la définition
des brains.

Les brains sont d'abord construits indépendamment. Les occurrences
de \texttt{bien} appartiennent à leurs brains respectifs, sans
être encore fusionnées.

\section{Fusion}

Les occurrences ayant l'identité canonique \texttt{bien} sont
rattachées au même nœud du maillage global.

Les directions naturelles observées comprennent notamment :
\[
\texttt{bien}
\overset{\mathrm{Nat}}{\longrightarrow}
\texttt{normalement}
\]
et
\[
\texttt{bien}
\overset{\mathrm{Nat}}{\longrightarrow}
\texttt{aujourd'hui}.
\]

Dans les autres trajectoires, les relations naturelles suivantes
sont également conservées :
\[
\texttt{bien}
\overset{\mathrm{Nat}}{\longrightarrow}
\texttt{de},
\qquad
\texttt{bien}
\overset{\mathrm{Nat}}{\longrightarrow}
\texttt{que}.
\]

Le nœud \texttt{bien} devient ainsi un point d'accès vers un
écosystème de directions issues de plusieurs trajectoires.

Les relations anti-causales déclarées dans le modèle sont
conservées séparément :
\[
\begin{aligned}
\texttt{bien}
&\overset{\mathrm{Anti}}{\longrightarrow}\texttt{c'est},\\
\texttt{bien}
&\overset{\mathrm{Anti}}{\longrightarrow}\texttt{oui},\\
\texttt{bien}
&\overset{\mathrm{Anti}}{\longrightarrow}\texttt{vais}.
\end{aligned}
\]

Ces relations ne sont pas automatiquement transformées en
continuations naturelles.

\section{Nouveau prompt : \texttt{va-t-il bien aujourd'hui ?}}

Le prompt courant est tokenisé selon la règle choisie. Pour
simplifier la notation, on représente ses ancres par :
\[
P=(\texttt{va},\texttt{il},\texttt{bien},
\texttt{aujourd'hui}).
\]

Le maillage contextuel est construit séparément du maillage
global. Il contient une colonne vertébrale ordonnée avec des
emplacements intermédiaires, ainsi que les relations
multidirectionnelles prévues par la règle de construction.

La recherche ne se contente pas de constater que le token
\texttt{bien} existe dans le maillage. Elle recherche un
alignement qui prend en compte la configuration du prompt :
\[
\{\texttt{va},\texttt{il},\texttt{bien},
\texttt{aujourd'hui}\}.
\]

Les correspondances avec \texttt{va}, \texttt{il} et
\texttt{aujourd'hui} peuvent contribuer à identifier des chemins
compatibles avec le prompt. Elles ne suffisent pas isolément :
il faut également une trajectoire admissible et une orientation
cohérente.

\section{Pourquoi l'écosystème importe}

Une fois \texttt{bien} retrouvé, la recherche examine les
directions et relations de son écosystème, mais ne les traite
pas toutes de façon identique.

La direction
\[
\texttt{bien}
\overset{\mathrm{Nat}}{\longrightarrow}
\texttt{aujourd'hui}
\]
est directement soutenue par une trajectoire naturelle
compatible avec certains éléments du prompt.

Les relations anti-causales vers \texttt{c'est}, \texttt{oui}
et \texttt{vais} conservent un statut différent. Elles peuvent
servir à distinguer ou réorienter l'exploration, mais ne doivent
pas être assimilées automatiquement à des continuations naturelles.

L'ensemble du prompt guide la sélection ; l'écosystème global
apporte les directions disponibles ; la trajectoire orientée
contraint la recherche.

\section{Limite de l'exemple}

Cet exemple illustre le mécanisme, mais ne suffit pas à déterminer
à lui seul une réponse exacte à la question. Il faudrait connaître
le dataset complet, les brains effectivement construits, les
relations enregistrées et la règle de traitement des relations
anti-causales.

La théorie doit éviter de fabriquer une continuation précise
lorsque le maillage disponible ne la soutient pas.

% ============================================================
\chapter{Théorie des empreintes : définitions et usage dans le modèle}
% ============================================================

\section{Pourquoi introduire les empreintes}

La théorie des empreintes fournit un cadre pour décrire ce qu'un
protocole d'observation distingue effectivement. Elle ne remplace
ni la construction des brains ni la fusion des identités.

Dans le présent modèle, elle sert à caractériser des tokens,
des trajectoires ou des alignements selon un ensemble déclaré
d'observations. Elle permet de préciser quand deux objets sont
indiscernables relativement à ce protocole.

Il faut distinguer :
\[
\text{identité lexicale}
\quad\neq\quad
\text{équivalence observationnelle}.
\]

La première règle décide quels tokens partagent une identité
canonique. La seconde décide quels objets donnent les mêmes
résultats sous un protocole d'observation.

\section{Espaces d'observation et sondes}

\begin{definition}[Objets observés, contextes et sondes]
Soient :
\begin{itemize}
  \item \(\X\), un ensemble d'objets observables ;
  \item \(\mathcal P\), un ensemble de contextes ;
  \item \(\F\), une famille de sondes ;
  \item \(\Phi\), une fonction d'observation.
\end{itemize}

On pose :
\[
\Phi:\mathcal P\times\X\times\F\longrightarrow\R.
\]

Pour un contexte \(P\), un objet \(x\) et une sonde \(f\),
\[
\Phi(P,x,f)
\]
est la valeur observée lorsque le protocole applique \(f\) à
l'objet \(x\) dans le contexte \(P\).
\end{definition}

Une sonde est une opération d'observation définie par le
protocole. Elle ne doit pas être présentée comme une source
automatique de vérité sémantique.

\section{Empreinte informationnelle}

\begin{definition}[Empreinte informationnelle]
Pour une famille finie de sondes
\[
\F_K=\{f_1,\ldots,f_K\},
\]
l'empreinte informationnelle d'un objet \(x\) dans le contexte
\(P\) est :
\[
I_{\Phi,\F_K}^{P}(x)
=
\bigl(
\Phi(P,x,f_1),\ldots,\Phi(P,x,f_K)
\bigr).
\]
\end{definition}

L'empreinte est donc une représentation produite par un protocole
d'observation déterminé. Elle ne représente pas une essence
absolue de l'objet, indépendamment de toute observation.

\section{Application aux objets du maillage}

Dans cette théorie, les objets observés peuvent notamment être :
\[
x\in
\{
\text{token},
\text{occurrence},
\text{brain},
\text{trajectoire},
\text{alignement},
\text{continuation candidate}
\}.
\]

Les sondes doivent être définies en fonction du type d'objet.
Par exemple, pour une trajectoire candidate, un protocole peut
observer :
\begin{itemize}
  \item si les ancres du prompt sont présentes dans l'ordre exigé ;
  \item si les relations intermédiaires sont admissibles ;
  \item si l'orientation de recherche est respectée ;
  \item quelles relations naturelles soutiennent la continuation ;
  \item quelles relations anti-causales sont activées ;
  \item si la provenance des relations est disponible ;
  \item si une condition de terminaison est atteinte.
\end{itemize}

Ces exemples définissent des familles de sondes possibles. Ils
ne signifient pas que leurs sorties numériques sont déjà fixées :
le protocole doit spécifier leur valeur et leurs conditions
d'application avant qu'elles soient utilisées.

\section{Indiscernabilité observationnelle}

\begin{definition}[Indiscernabilité]
Pour un contexte \(P\) et une famille de sondes \(\F_K\), on
définit :
\[
x\sim_{\Phi,P,\F_K}y
\iff
I_{\Phi,\F_K}^{P}(x)
=
I_{\Phi,\F_K}^{P}(y).
\]
Deux objets sont indiscernables lorsque le protocole choisi
produit exactement la même empreinte pour les deux objets.
\end{definition}

\begin{proposition}
La relation \(\sim_{\Phi,P,\F_K}\) est une relation d'équivalence.
\end{proposition}

\begin{proof}
La réflexivité découle de
\(I(x)=I(x)\). La symétrie découle de
\(I(x)=I(y)\Rightarrow I(y)=I(x)\). La transitivité découle de
\(I(x)=I(y)\) et \(I(y)=I(z)\Rightarrow I(x)=I(z)\).
\end{proof}

\section{Fibres observationnelles et quotient}

\begin{definition}[Fibre observationnelle]
La fibre de \(x\) est la classe :
\[
[x]_{\Phi,P,\F_K}
=
\{y\in\X:y\sim_{\Phi,P,\F_K}x\}.
\]
Elle regroupe tous les objets indiscernables de \(x\) sous le
protocole déclaré.
\end{definition}

\begin{definition}[Quotient observationnel]
Le quotient est l'ensemble des classes d'indiscernabilité :
\[
\X/{\sim_{\Phi,P,\F_K}}.
\]
La projection canonique est :
\[
q(x)=[x]_{\Phi,P,\F_K}.
\]
\end{definition}

Le quotient décrit les distinctions que le protocole permet
effectivement de faire. Deux objets appartenant à la même fibre
ne peuvent pas être distingués par les sondes retenues, même
si leurs identités brutes sont différentes.

\section{Utilisation dans la théorie du maillage}

Les empreintes interviennent dans le présent modèle pour trois
usages.

\subsection{Comparer des trajectoires candidates}

On peut observer deux trajectoires candidates sous un même
protocole. Si leurs empreintes sont égales, elles sont
indiscernables relativement à ce protocole.

Cela ne prouve pas qu'elles sont identiques à tous les niveaux :
cela signifie uniquement que les sondes choisies ne les
distinguent pas.

\subsection{Comparer des alignements}

Deux alignements différents peuvent satisfaire les mêmes
observables de correspondance. Ils peuvent alors appartenir à
la même fibre observationnelle.

Cette situation indique que le protocole actuel ne permet pas
de les distinguer. Elle ne justifie pas d'inventer une différence
qui n'est pas observée.

\subsection{Évaluer le raffinement d'un protocole}

Si une famille de sondes \(\F_1\) est incluse dans une famille
\(\F_2\), alors :
\[
\F_1\subseteq\F_2
\]
et
\[
I_{\Phi,\F_2}^{P}(x)=I_{\Phi,\F_2}^{P}(y)
\Rightarrow
I_{\Phi,\F_1}^{P}(x)=I_{\Phi,\F_1}^{P}(y).
\]

Une famille plus riche peut distinguer davantage d'objets, mais
elle ne les distingue effectivement que si ses nouvelles sondes
produisent des observations différentes.

\section{Pourquoi l'empreinte ne remplace pas la fusion lexicale}

\begin{proposition}[Distinction entre fusion lexicale et quotient
observationnel]
La fusion de deux occurrences de même identité canonique ne
découle pas automatiquement de leur indiscernabilité
observationnelle. Réciproquement, deux objets indiscernables
sous un protocole ne doivent pas nécessairement être fusionnés
en une seule identité lexicale.
\end{proposition}

\begin{proof}
La fusion lexicale est définie par une règle d'identité portant
sur les tokens. L'indiscernabilité est définie par l'égalité
des sorties d'un protocole d'observation. Ces deux relations
ont des définitions et des objectifs distincts ; aucune
implication générale entre elles n'est donc établie sans
hypothèse supplémentaire.
\end{proof}

Cette distinction est importante pour éviter de confondre
l'identité du token avec son rôle dans une trajectoire donnée.

\section{Séparation entre observation et sémantique}

Une empreinte est un résultat d'observation. Une interprétation
sémantique de cette empreinte nécessite un protocole d'interprétation
supplémentaire.

De même, une relation du maillage peut être observée et enregistrée
sans que sa signification sémantique soit automatiquement
certifiée.

La théorie du maillage utilise donc les empreintes pour rendre
les distinctions et les comparaisons explicites, et non pour
transformer une représentation en vérité par simple définition.

% ============================================================
\chapter{Propriétés établies, hypothèses et limites}
% ============================================================

\section{Propriétés établies par construction}

Les propriétés suivantes découlent des définitions, sous réserve
que les structures requises soient effectivement conservées.

\begin{enumerate}[label=\arabic*.]
  \item Chaque exemple peut être représenté par un brain séparé.
  \item Les occurrences de tokens peuvent partager une identité
        canonique dans le maillage global.
  \item La fusion des identités est compatible avec la conservation
        des trajectoires, si les relations et leur provenance
        sont conservées.
  \item Le maillage contextuel peut être construit séparément du
        maillage global.
  \item Les variables contextuelles peuvent accepter des
        structures intermédiaires, si leur règle d'affectation
        l'autorise.
  \item Un alignement exige plus que la simple présence des tokens :
        il exige des relations et une orientation admissibles.
  \item Une réponse sans relation de soutien vers un alignement
        du prompt ne satisfait pas la définition de recevabilité
        structurelle.
  \item Les relations naturelles et anti-causales peuvent être
        distinguées par leurs types.
  \item Les empreintes définissent une indiscernabilité
        observationnelle et les fibres associées.
\end{enumerate}

\section{Hypothèses de fonctionnement à tester}

Les affirmations suivantes sont des hypothèses de conception
ou des hypothèses expérimentales, pas des théorèmes issus des
seules définitions :

\begin{enumerate}[label=\arabic*.]
  \item Le maillage global contient suffisamment de trajectoires
        pertinentes pour répondre à une proportion utile de
        prompts nouveaux.
  \item Les alignements contextuels éliminent une partie importante
        des continuations sans rapport avec la question.
  \item La conservation des directions naturelles améliore la
        sélection des continuations.
  \item Le traitement des relations anti-causales réduit les
        continuations structurellement incompatibles.
  \item La combinaison du prompt complet, des chemins activés et
        de l'historique limite les boucles de génération.
  \item La navigation dans le maillage permet de produire des
        réponses nouvelles qui ne sont pas seulement la copie
        d'une réponse unique du dataset.
\end{enumerate}

Ces hypothèses doivent être évaluées par des expériences où les
conditions de test, les métriques et les cas d'échec sont définis.

\section{Cohérence structurelle et vérité factuelle}

\begin{proposition}[Absence de garantie factuelle par la seule
cohérence relationnelle]
La cohérence structurelle d'une réponse dans le maillage global
ne suffit pas à démontrer sa vérité factuelle.
\end{proposition}

\begin{proof}
Le maillage est construit à partir de données et de règles
d'étiquetage. Si les données contiennent une affirmation fausse
et que les relations correspondantes sont conservées, la
structure peut rester cohérente avec cette affirmation fausse.
La cohérence interne ne suffit donc pas à établir la vérité
externe.
\end{proof}

Cette limite n'annule pas le mécanisme de recherche. Elle
précise ce qu'il permet d'établir : une réponse structurellement
soutenue, non une vérité garantie indépendamment de ses sources.

\section{Limites opérationnelles et traitement}

\subsection{Ambiguïté lexicale}

Un même token canonique peut apparaître dans plusieurs contextes
et avoir des rôles différents. La théorie conserve ces différences
par les trajectoires, les relations et les empreintes éventuelles.
Elle ne doit pas supposer que l'identité lexicale implique une
identité de rôle.

\subsection{Relations anti-causales insuffisamment définies}

Si l'étiquette anti-causale n'a aucune règle opérationnelle,
le système ne peut pas l'utiliser de manière reproductible.

Le contournement correct n'est pas de supprimer cette catégorie,
mais de déclarer une règle de classification, une provenance et
un effet sur la recherche. Les relations non classées doivent
rester non classées plutôt que recevoir une étiquette inventée.

\subsection{Absence d'alignement}

Si
\[
\Align(P,G)=\varnothing,
\]
aucun alignement admissible n'a été trouvé selon le protocole
courant.

Le système doit signaler cette absence, élargir la recherche
selon une règle explicite, ou demander de nouvelles données.
Il ne doit pas présenter une continuation arbitraire comme
si elle avait été soutenue par le maillage.

\subsection{Plusieurs alignements compatibles}

Plusieurs alignements peuvent satisfaire toutes les contraintes.
La théorie n'impose pas de prétendre qu'un seul est objectivement
meilleur lorsque les observations ne les distinguent pas.

Le système peut conserver plusieurs branches, présenter plusieurs
réponses candidates ou utiliser une politique de sortie
déterministe clairement déclarée. Cette dernière est une règle
d'implémentation, pas une conséquence automatique de la fusion.

\subsection{Coût de recherche}

Un maillage peut être très grand et contenir de nombreux chemins.
La théorie définit les contraintes de pertinence, mais ne garantit
pas que leur recherche exhaustive soit toujours peu coûteuse.

Une implémentation peut utiliser des index, des caches ou des
bornes de recherche, à condition que ces optimisations ne soient
pas présentées comme des preuves de validité. Toute restriction
de recherche doit être documentée, car elle peut exclure des
alignements qui existent dans le maillage complet.

% ============================================================
\chapter{Protocole de validation expérimentale}
% ============================================================

\section{Objectif}

L'objectif expérimental est de déterminer si le mécanisme
\[
\text{brains}
\to
\text{fusion}
\to
\text{alignement}
\to
\text{navigation}
\]
produit effectivement des réponses plus pertinentes et plus
structurellement soutenues que des mécanismes de référence.

L'expérience ne doit pas confondre :
\begin{itemize}
  \item l'existence d'un maillage ;
  \item la possibilité de trouver un alignement ;
  \item la qualité de la réponse produite ;
  \item la vérité factuelle de la réponse ;
  \item le coût de calcul.
\end{itemize}

\section{Conditions à comparer}

On peut comparer plusieurs configurations :
\begin{enumerate}[label=\arabic*.]
  \item trajectoires sans fusion d'identités ;
  \item maillage avec fusion, mais sans maillage contextuel ;
  \item maillage avec alignement contextuel, sans contraintes
        anti-causales ;
  \item maillage avec alignement, orientation et relations
        naturelles ;
  \item modèle complet avec relations naturelles et
        anti-causales, historique de génération et gestion des
        boucles.
\end{enumerate}

Chaque configuration doit utiliser le même dataset et des
conditions d'évaluation comparables.

\section{Mesures à enregistrer}

Les mesures possibles comprennent :
\begin{itemize}
  \item proportion de prompts possédant au moins un alignement ;
  \item proportion de réponses disposant d'une provenance
        structurelle traçable ;
  \item taux de réponses sans relation admissible vers le prompt ;
  \item taux de boucles sans progression ;
  \item nombre d'alignements compatibles par prompt ;
  \item coût de recherche par prompt ;
  \item qualité de pertinence évaluée selon un protocole déclaré ;
  \item exactitude factuelle, lorsque des références vérifiables
        sont disponibles ;
  \item proportion de réponses nouvelles par rapport aux réponses
        brutes du dataset.
\end{itemize}

Les métriques de qualité doivent être définies avant les tests.
Aucun indicateur isolé ne prouve à lui seul la validité complète
de la théorie.

\section{Provenance et auditabilité}

Chaque réponse candidate devrait pouvoir être accompagnée d'un
certificat interne indiquant :
\begin{itemize}
  \item les ancres du prompt retrouvées ;
  \item les valeurs affectées aux variables \(c_i\) ;
  \item les chemins utilisés ;
  \item les relations naturelles prises en compte ;
  \item les relations anti-causales examinées ;
  \item les brains ou exemples sources pertinents ;
  \item les règles de recherche appliquées ;
  \item les éventuelles contraintes de calcul ou de troncature.
\end{itemize}

Ce certificat ne prouve pas la vérité de la réponse. Il permet
de vérifier si la réponse est effectivement soutenue par le
mécanisme annoncé.

% ============================================================
\chapter{Synthèse formelle}
% ============================================================

\section{Architecture complète}

Le modèle peut être résumé par les objets suivants :
\[
\begin{aligned}
\D&=\text{dataset de couples prompt/réponse},\\
B_i&=\text{brain de l'exemple }i,\\
G&=\text{maillage global fusionné},\\
C_P&=\text{maillage contextuel du prompt }P,\\
\Align(P,G)&=\text{ensemble des alignements admissibles},\\
\A(P,G)&=\text{région activée par ces alignements},\\
S_t&=\text{état de recherche à l'étape }t,\\
H_t&=\text{historique de génération}.
\end{aligned}
\]

La construction hors ligne est :
\[
\boxed{
\D
\longrightarrow
\{B_i\}_{i=1}^{n}
\longrightarrow
G
}
\]

La recherche en ligne est :
\[
\boxed{
P
\longrightarrow
C_P
\longrightarrow
\Align(P,G)
\longrightarrow
\A(P,G)
\longrightarrow
S_t
\longrightarrow
R
}
\]

La relation entre \(C_P\) et \(G\) est une relation de recherche,
pas une fusion :
\[
\boxed{C_P\text{ interroge }G,\quad C_P\text{ ne s'incorpore pas à }G.}
\]

La cohérence structurelle est définie par :
\[
\boxed{
\begin{gathered}
\text{alignement contextuel}\\
+\ \text{orientation }c_0\to c_x\\
+\ \text{compatibilité avec les directions naturelles}\\
+\ \text{traitement des relations anti-causales}\\
+\ \text{traçabilité des relations de soutien}.
\end{gathered}
}
\]

\section{Formulation finale de la théorie}

La théorie construit, à partir d'un dataset, un brain distinct
pour chaque exemple prompt/réponse. Chaque brain conserve les
occurrences de tokens, l'ordre des trajectoires, les relations
observées et leur provenance. Les identités communes sont
ensuite fusionnées dans un maillage global : les tokens deviennent
des points d'accès partagés, tandis que les relations issues des
trajectoires restent conservées.

Chaque token canonique possède un écosystème relationnel. Cet
écosystème distingue les directions naturelles et les relations
anti-causales, dont le rôle doit être déclaré et appliqué pendant
la recherche. La fusion factorise les identités communes sans
réduire l'ensemble relationnel à un voisinage indifférencié.

Lorsqu'un prompt arrive, le système construit un maillage
contextuel indépendant. Les tokens du prompt y forment des ancres
reliées par une structure multidirectionnelle. Les variables
\(c_i\) acceptent des tokens, des ensembles ou des trajectoires
intermédiaires du maillage global. Elles rendent ainsi possible
l'alignement de plusieurs ancres sans exiger qu'elles soient
directement voisines.

Le parcours de recherche est orienté par \(c_0\to c_x\). La
présence des ancres ne suffit pas : un alignement admissible
doit respecter la structure contextuelle et son orientation.
Les régions activées par les alignements servent ensuite à
rechercher des continuations issues des trajectoires du dataset
et des relations du maillage global.

La réponse en construction reste soumise au prompt complet,
aux éléments déjà produits, aux directions naturelles et au
traitement des relations anti-causales. La recherche n'est donc
pas réduite au choix local du voisin le plus fréquent : elle
navigue dans un maillage de trajectoires factorisées sous des
contraintes contextuelles et relationnelles.

La théorie des empreintes complète cette architecture en
définissant comment les tokens, trajectoires et alignements
peuvent être observés, comparés et regroupés selon leurs
empreintes informationnelles. Elle ne remplace ni la règle
d'identité lexicale ni les contraintes de recherche. Elle rend
explicites les distinctions effectivement visibles sous un
protocole d'observation.

Le résultat visé est une génération dont les continuations
doivent être soutenues par un alignement et une trajectoire
compatibles avec l'organisation du maillage. La théorie définit
ainsi un mécanisme interne de contrainte et de discrimination
des réponses candidates, sans confondre cohérence structurelle,
validité sémantique et vérité factuelle.

\end{document}